% GENERATED FILE. Edit canonical lessons, then run npm run generate:books.
% canonical-source-hash: fnv1a64:ea387e6758bd8bf1
% canonical-lessons: GU-C133-kathan, GU-C133-sahelu, GU-C133-agharu, GU-C133-hosiyar, GU-C133-alsu

\chapter{Hard, Easy, Difficult, Clever, Lazy}
\label{ch:gu-hard-easy-difficult-clever-lazy}
\hlchaptermodality{133}

\begin{chapteropening}
\textbf{By the end of this chapter:} \emph{I can say hard, easy, difficult, clever and lazy.}

Complete the last lesson of chapter 133: i can say hard, easy, difficult, clever and lazy.
\end{chapteropening}

\section[kaṭhaṇ]{\gu{કઠણ} (kaṭhaṇ) — hard}

\label{lesson:GU-C133-kathan}



\begin{quote}
Before the new one: say the Gujarati for spicy, then the Gujarati for soft.
\end{quote}

\subsection*{You'll want to know: \gu{કઠણ}}
\textbf{\gu{કઠણ}} (\emph{kaṭhaṇ}) — \textquotedblleft{}hard\textquotedblright{}.

\begin{grammarlens}[title={What you've built}]
One more word for this chapter.
\end{grammarlens}

\subsection*{Your turn}
\begin{itemize}
\raggedright
  \item \emph{Say it:} \emph{kaṭhaṇ}
  \item \emph{Say it:} \emph{kaṭhaṇ}, once more
  \item \emph{Say it:} say \emph{kaṭhaṇ}
  \item \emph{Recall:} say the Gujarati for spicy, then the Gujarati for soft, then say \emph{kaṭhaṇ} again
\end{itemize}

\begin{tcolorbox}[breakable,colback=teal!4,colframe=teal!35!black,title={Before you move on}]
What does \gu{કઠણ} mean? (\textquotedblleft{}Hard\textquotedblright{}.) Say it once more.
\end{tcolorbox}
\section[sahelũ]{\gu{સહેલું} (sahelũ) — easy}

\label{lesson:GU-C133-sahelu}



\begin{quote}
Before the new one: say the Gujarati for soft, then the Gujarati for hard.
\end{quote}

\subsection*{You'll want to know: \gu{સહેલું}}
\textbf{\gu{સહેલું}} (\emph{sahelũ}) — \textquotedblleft{}easy\textquotedblright{}.

\begin{grammarlens}[title={What you've built}]
One more word for this chapter.
\end{grammarlens}

\subsection*{Your turn}
\begin{itemize}
\raggedright
  \item \emph{Say it:} \emph{sahelũ}
  \item \emph{Say it:} \emph{sahelũ}, once more
  \item \emph{Say it:} say \emph{sahelũ}
  \item \emph{Recall:} say the Gujarati for soft, then the Gujarati for hard, then say \emph{sahelũ} again
\end{itemize}

\begin{tcolorbox}[breakable,colback=teal!4,colframe=teal!35!black,title={Before you move on}]
What does \gu{સહેલું} mean? (\textquotedblleft{}Easy\textquotedblright{}.) Say it once more.
\end{tcolorbox}
\section[agharũ]{\gu{અઘરું} (agharũ) — difficult}

\label{lesson:GU-C133-agharu}



\begin{quote}
Before the new one: say the Gujarati for hard, then the Gujarati for easy.
\end{quote}

\subsection*{You'll want to know: \gu{અઘરું}}
\textbf{\gu{અઘરું}} (\emph{agharũ}) — \textquotedblleft{}difficult\textquotedblright{}.

\begin{grammarlens}[title={What you've built}]
One more word for this chapter.
\end{grammarlens}

\subsection*{Your turn}
\begin{itemize}
\raggedright
  \item \emph{Say it:} \emph{agharũ}
  \item \emph{Say it:} \emph{agharũ}, once more
  \item \emph{Say it:} say \emph{agharũ}
  \item \emph{Recall:} say the Gujarati for hard, then the Gujarati for easy, then say \emph{agharũ} again
\end{itemize}

\begin{tcolorbox}[breakable,colback=teal!4,colframe=teal!35!black,title={Before you move on}]
What does \gu{અઘરું} mean? (\textquotedblleft{}Difficult\textquotedblright{}.) Say it once more.
\end{tcolorbox}
\section[hośiyār]{\gu{હોશિયાર} (hośiyār) — clever}

\label{lesson:GU-C133-hosiyar}



\begin{quote}
Before the new one: say the Gujarati for easy, then the Gujarati for difficult.
\end{quote}

\subsection*{You'll want to know: \gu{હોશિયાર}}
\textbf{\gu{હોશિયાર}} (\emph{hośiyār}) — \textquotedblleft{}clever\textquotedblright{}.

\begin{grammarlens}[title={What you've built}]
One more word for this chapter.
\end{grammarlens}

\subsection*{Your turn}
\begin{itemize}
\raggedright
  \item \emph{Say it:} \emph{hośiyār}
  \item \emph{Say it:} \emph{hośiyār}, once more
  \item \emph{Say it:} say \emph{hośiyār}
  \item \emph{Recall:} say the Gujarati for easy, then the Gujarati for difficult, then say \emph{hośiyār} again
\end{itemize}

\begin{tcolorbox}[breakable,colback=teal!4,colframe=teal!35!black,title={Before you move on}]
What does \gu{હોશિયાર} mean? (\textquotedblleft{}Clever\textquotedblright{}.) Say it once more.
\end{tcolorbox}
\section[āḷsu]{\gu{આળસુ} (āḷsu) — lazy}

\label{lesson:GU-C133-alsu}



\begin{quote}
Before the new one: say the Gujarati for difficult, then the Gujarati for clever.
\end{quote}

\subsection*{You'll want to know: \gu{આળસુ}}
\textbf{\gu{આળસુ}} (\emph{āḷsu}) — \textquotedblleft{}lazy\textquotedblright{}.

\begin{grammarlens}[title={What you've built}]
One more word for this chapter.
\end{grammarlens}

\subsection*{Your turn}
\begin{itemize}
\raggedright
  \item \emph{Say it:} \emph{āḷsu}
  \item \emph{Say it:} \emph{āḷsu}, once more
  \item \emph{Say it:} say \emph{āḷsu}
  \item \emph{Recall:} say the Gujarati for difficult, then the Gujarati for clever, then say \emph{āḷsu} again
\end{itemize}

\begin{tcolorbox}[breakable,colback=teal!4,colframe=teal!35!black,title={Before you move on}]
What does \gu{આળસુ} mean? (\textquotedblleft{}Lazy\textquotedblright{}.) Say it once more.
\end{tcolorbox}
